%\vspace{-1em}
\section{Related Work}
%\vspace{-1em}
Previous research has witnessed a significant focus on privacy protection in the mobile area. 
These works highlight the intricate privacy dimensions specific to mobile apps. Specifically, these studies offer valuable perspectives on several key areas: they explore the privacy challenges in super-app ecosystems~\cite{dongexploring, 10.1145/3576915.3616676, 10.1145/3576915.3616591, 10.5555/3620237.3620608, yang2022cross, zhang2022identity, haney2021s, chen2019demystifying}, investigate methods through which data might be stealthily exfiltrated or mishandled~\cite{nan2023are, koch2023ok,diamantaris2021sneaky}, assess the alignment between the mobile app and regulations~\cite{pantrap, khandelwal2023unpacking, 10.1145/3576915.3623067, xiao2023lalaine, meng2023post, li2022collect, yang2022cross, young2022skilldetective, balash2022security, bui2021consistency, nguyen2021share, zuo2019does}, and propose strategies to enhance privacy compliance and transparency within apps~\cite{klein2023general,ferreira2023rulekeeper,jordan2021viceroy,nguyen2022freely,wang2022srr,cai2024famos}.
% 
These studies emphasize the complexity of privacy issues, covering personal and device information, location, camera, microphone, etc. These investigations reflect the community's commitment to tackling mobile app privacy and security challenges. %


However, an unexplored area remains regarding the potential privacy risks in the mini-app interaction history within super-app ecosystems. Unlike location or microphone data, these data do not typically trigger system-level permission prompts, making them less visible and hence less concerned. Meanwhile, they can reveal extensive insights into user privacy attributes. In this paper, we close a critical gap in understanding how seemingly non-sensitive mini-app interaction history can pose novel risks, underscoring the need for more nuanced privacy protections for the super-app ecosystems.
